%% 08_lna.tex -- Chapter 8: The AC-Coupled Biopotential LNA
%% Status: OUTLINE (build stage 3). Replace the outline with the full chapter when it is written.
\chapter{The AC-Coupled Biopotential LNA}\label{ch:lna}
\chapterlead{Two capacitors, two pseudo-resistors and an OTA: the Harrison amplifier is the circuit every neural and ExG front-end is compared with.}

\begin{specnote}[Chapter status: outline]
This chapter is an outline. It will be written in build stage 3 (Part II with Labs 3--5). The topics, the design question and the labs below are fixed; the derivations, worked examples and figures are still to come.
\end{specnote}

\begin{objectives}
\begin{itemize}
  \item derive the gain, band edges and noise of the capacitive-feedback LNA;
  \item explain how MOS-bipolar pseudo-resistors reach high-pass corners in the millihertz range;
  \item design the OTA (folded cascode or current-scaled telescopic) for a given noise and bandwidth;
  \item budget gain, bandwidth and noise for ECG, EEG and EMG modes;
  \item recognise the linearity, leakage and PVT problems of pseudo-resistors.
\end{itemize}
\end{objectives}

\begin{insight}[Design question]
How do you size the capacitors, pseudo-resistors and OTA so that one AC-coupled LNA meets the noise, bandwidth and power targets of every \OB{} mode?
\end{insight}

\section{The Harrison topology}\label{sec:ch8-1}
Midband gain $A_\mathrm{M} = C_1/C_2$, low corner $f_\mathrm{L} = 1/(2\pi R_\mathrm{pr} C_2)$ and high corner set by the OTA transconductance and load; the original design reaches \qtyrange{0.025}{7200}{\hertz} with \qty{2.2}{\uVrms} of noise at \qty{80}{\micro\watt}~\cite{harrison2003}.

\section{Pseudo-resistors}\label{sec:ch8-2}
Diode-connected MOS-bipolar devices give teraohm small-signal resistance in a tiny area; their resistance depends strongly on voltage, temperature and leakage, so the high-pass corner moves.

\section{OTA design}\label{sec:ch8-3}
Folded cascode versus current-scaled telescopic OTAs; putting most of the current into the input pair to approach the NEF limit of \cref{ch:noise}; capacitive noise gain $(C_1+C_2+C_\mathrm{in})/C_2$.

\section{Gain, bandwidth and noise budget}\label{sec:ch8-4}
Allocating the \qty{1.5}{\uVrms} EEG-mode noise target, choosing $C_1$ for the input impedance it presents (\cref{ch:electrodes}), and switching bandwidth between modes.

\section{Linearity and pseudo-resistor problems}\label{sec:ch8-5}
Distortion at large EMG and artefact amplitudes, leakage at high temperature, and tunable alternatives (bias-controlled pseudo-resistors, duty-cycled resistors).

\begin{tryit}[Labs for this chapter]
Lab~4 (\cref{sec:lab4}) designs the capacitive-feedback LNA with pseudo-resistors in \sgx{}: AC, noise and transient simulations and the NEF versus bias current.
\end{tryit}

\paragraph{Read with this chapter} Harrison and Charles~\cite{harrison2003}, Yazicioglu, Van Hoof and Puers~\cite{yazicioglubook}.
